% Transcription — fonds Alexandre Grothendieck, shelfmark 11, batch 3 (pages 41-57).
% Established from the facsimile archives/batches/11/batch-03.pdf (a mirror of
% https://grothendieck.umontpellier.fr/11.pdf), per the skill
% transcribe-grothendieck. Its PDF page 1 is archivists' page 41 (checked
% against the pencilled number); the batch file has no cover sheet.
%
% Pass: Opus 5 (claude-opus-5), 2026-09-14 — first pass, unchecked against the pages by a human.
%
% Revised 2026-09-23 (Opus 5.5), after the find-novelty pass.
%   - Header: pages 50ff were summarised « for X with H*(X) = k ⊕ k·y »; the
%     page states no such hypothesis (checked on the facsimile), and under it
%     the page's five-element « base » of H*(X×X) would be dependent. Now
%     described as the tautological part the heading names.
%   - Header: the process sentence made consistent with batch 2's (concurrent
%     passes; no batch fixed another's notation).
%   - Page 50, degree 2n « base » of H*(X×X): the page writes y_1y_2y_3
%     (re-read on the facsimile at 350 dpi), silently normalised to y_1y_2 by
%     the first pass; now transcribed as written, with a sic note.
%
% The last seventeen leaves of the folder. Eight carry his hand: 41, 43, 45,
% 47, 50, 52, 54, 56. The others are skipped: 42, 44, 46, 48, 51, 53, 55 and
% 57 are the backs of an English typescript that has nothing to do with these
% notes, and 49 has no writing of its own. The three batches of the folder were
% transcribed on the same day by concurrent passes, so no batch fixed the
% notation of another (batch 2's header says the same); page 39 was glanced
% at, page 40 is a typescript verso. No run crosses the 40/41
% boundary: page 41 opens on a fresh heading. Page 39's Proposition (subalgebras
% A_n ⊂ H*(X^n) stable under inverse and direct images by simplicial maps) is
% taken up again by the a)/b)/c) checklist of page 56.
%
% Pages 41-47, bivariant graded algebras and the formalism of correspondences.
% Page 41 defines the category: graded k-algebras with a degree d(A), zero
% outside [1, d(A)] as the page writes it, with κ_A : A → k; « jaugées » when
% equipped with a map to k. For simplicial ones, A(I' ⊔ I'') ≃ A(I') ⊗ A(I'')
% and (α ⊔ β)_* = α_* ⊗ β_*. Evaluation at Δ⁰ is an equivalence onto Poincaré
% algebras, and A(·) ↦ (A(0), Δ) is fully faithful into gauged algebras with a
% Δ ∈ [A⊗A]^{d(A)}. Example 1: a cohomology theory with Künneth and Poincaré
% applied to I ↦ X^I. Page 43: variants (compact topological manifolds, ring
% coefficients), example 2 (Chow ring, quotients by numerical, homological,
% τ-algebraic equivalence, Picard, Albanese), example 3 (K(X) gauged by χ).
% Then « Formalisme des correspondances »: C ∈ A(I ⊔ I') acts A(I) → A(I')
% through α^*, multiplication by C and α'_*, and the side of the product is a
% choice. Page 45 defines the composition v∘u = p_{31*}(p_{21}^*(u) p_{32}^*(v))
% and checks (v∘u)~ = ṽ ∘ ũ by projection formula and base change (« axiome
% cartésien oublié ! », « OK mod signes »). Page 47 defines Δ_I = δ_*(1) and
% checks that its action is the identity, then breaks off mid-sentence on
% associativity and units.
%
% Pages 50-56, « Sous-algèbre tautologique de H*(X^n) ». The page puts no
% hypothesis on X. Under that heading it lists, in items headed H*(X), H*(X×X),
% H*(X³): for X the class y ∈ H^n(X), y² = 0, « base » 1, y; then generators,
% relations and « bases » for X×X (y_1, y_2, δ, δ² = χ y_1y_2) and for X³ (the
% δ_{ij} and δ). What is listed is the part generated by y and the diagonals,
% not the whole of H*: were H*(X) = k ⊕ k·y, the five-element « base » of
% H*(X×X) would lie in a space of dimension 4. Then
% generators of H*(X^k) as transforms of y⊗1⊗⋯⊗1 and δ_r⊗1⊗⋯⊗1. Page 54
% indexes the classes by pairs (R, J), R a partition of I and J ∈ R, with
% Z_{R,J} = {x_i = a on J, x_i = x_j on each A ∈ R}, and works out
% ε_{R,J} ε_{R',J'} from the partition R'' generated by R, R'. Page 56 counts
% codimensions to get s + s' ≤ k + 1 (equality gives the diagonal δ_k, k = s+s'
% is left with a « ? »), sketches linear-algebra diamonds with codimensions
% p, q, r, checks stability under ⊠, inverse and direct images, and boxes the
% conclusion: generated via ⊠ by y and the δ_I.
%
% Notation fixed for the batch: ⊔ as \sqcup; A(I) the value on a finite set I;
% X^I; lower and upper stars for push-forward and pull-back; the projections
% p_{ij}, p_i; ũ for the operator of a correspondence u; δ for a diagonal,
% Δ_I for its class; κ_A the gauge; χ the Euler characteristic, as he says;
% ⊠ as \boxtimes; 𝔓(I) as \mathfrak{P}(I). His interlinear insertions are
% set in the line where his caret puts them; \add is kept for editorial
% additions only.
\documentclass[11pt,a4paper]{article}
% Le préambule commun à toutes les transcriptions du fonds Grothendieck.
%
% Il tient en une idée : **l'incertitude se compose, elle ne se résout pas.**
% Une transcription qui lisse un mot douteux a détruit la seule chose pour
% laquelle on la faisait — un lecteur doit pouvoir savoir, à chaque endroit,
% ce qui était sur le feuillet et ce qui a été ajouté. Les six macros qui
% suivent sont donc l'appareil critique, et non de la décoration.
%
% Les mêmes commandes sont reconnues par `scripts/render.mjs`, qui produit la
% vue de lecture du site. Une macro ajoutée ici doit l'être aussi là-bas,
% faute de quoi le rendu échouera bruyamment — ce qui est voulu : un rendu
% silencieusement faux est pire qu'un rendu absent.

\ProvidesPackage{grothendieck}[2026/08/08 Transcriptions du fonds Grothendieck]

\usepackage{amsmath,amssymb,amsthm}
\usepackage{mathrsfs}
\usepackage{stmaryrd}
% Les diagrammes commutatifs sont partout dans ce fonds : c'est la langue
% même de la géométrie algébrique de Grothendieck, et une transcription qui
% les rendrait en prose ne transcrirait rien.
\usepackage{tikz-cd}
% Français par défaut — c'est la langue des feuillets — mais l'anglais est
% chargé aussi : la traduction et le résumé sont des documents anglais, et la
% typographie française (espaces avant les deux-points) y serait une faute.
% Ces documents basculent par \selectlanguage{english} en tête de corps.
\usepackage[english,french]{babel}
\usepackage{fontspec}
\usepackage{geometry}
\geometry{a4paper,margin=25mm,marginparwidth=28mm}
\usepackage{marginnote}
\usepackage{xcolor}
% ulem, not soul. `soul`'s \st and \ul are built on a character-by-character
% scan that XeLaTeX's font machinery breaks: the document fails with an
% undefined control sequence rather than degrading. ulem does the same two jobs
% and is Unicode-engine safe. `normalem` stops it hijacking \emph, which the
% transcriptions use for Grothendieck's own emphasis.
\usepackage[normalem]{ulem}
\usepackage{hyperref}
\hypersetup{colorlinks=true,linkcolor=black,urlcolor=black,pdfborder={0 0 0}}

% --- Métadonnées du lot -----------------------------------------------------
% Reprises telles quelles par le script de rendu, qui les affiche en tête de
% la vue de lecture. Elles fixent ce que le fichier prétend couvrir : sans
% elles, une transcription flotte sans point d'attache dans un fonds de seize
% mille feuillets.

\newcommand{\folder}[1]{\gdef\grothFolder{#1}}
\newcommand{\batch}[1]{\gdef\grothBatch{#1}}
\newcommand{\leaves}[2]{\gdef\grothFirst{#1}\gdef\grothLast{#2}}
\newcommand{\foldertitle}[1]{\gdef\grothFoldertitle{#1}}
\gdef\grothFolder{?}\gdef\grothBatch{?}\gdef\grothFirst{?}\gdef\grothLast{?}\gdef\grothFoldertitle{}

% --- Le résumé --------------------------------------------------------------
%
% Il ouvre la lecture modernisée, sous le titre « Résumé », et il est composé
% en petit. Ce n'est pas un ornement : c'est le seul passage du document qui
% ne soit pas des mathématiques, et un lecteur doit pouvoir le reconnaître
% avant de l'avoir lu — puis le sauter s'il connaît déjà le sujet, ou s'y
% arrêter s'il ne le connaît pas. Le filet qui le suit marque la reprise du
% corps.

\newenvironment{resume}
  {\small\setlength{\parskip}{0.45em}}
  {\par\medskip\hrule\medskip}

% --- Mots-clés ---------------------------------------------------------------
%
% En anglais, en clôture du résumé de la lecture modernisée : le vocabulaire
% moderne sous lequel chercher ce que les feuillets construisent. Le manifeste
% du site (`npm run manifest`) les extrait pour étiqueter la cote entière —
% cette ligne est la source unique des tags, il n'y a pas de fichier à côté.
\newcommand{\keywords}[1]{\par\smallskip\noindent
  {\footnotesize\textsc{Keywords} --- \textit{#1}}\par}

% --- L'appareil critique ----------------------------------------------------

% Le numéro de feuillet, en marge. C'est l'ancre qui permet de régler un
% désaccord sur une formule en regardant *un* feuillet plutôt qu'une chemise
% de six cents. Le site s'en sert aussi pour tourner les pages du fac-similé
% quand on fait défiler la transcription.
\newcommand{\leaf}[1]{%
  \marginnote{\footnotesize\color{gray!70}\textsf{#1}}%
  \label{leaf:#1}\ignorespaces}

% Illisible : on ne devine pas. Un mot inventé qui se lit comme les autres est
% le pire résultat possible.
\newcommand{\ill}{\textcolor{red!60!black}{[\,\ldots\,]}}

% Lecture douteuse : le mot est donné, et signalé comme incertain.
\newcommand{\uncertain}[1]{\uline{#1}}

% Ajout de l'éditeur — une lettre manquante, un mot sous-entendu.
\newcommand{\add}[1]{\textcolor{blue!50!black}{[#1]}}

% Biffé par Grothendieck lui-même. Ce qu'il a rayé fait partie du document :
% une rature dit souvent où la pensée a changé de direction.
\newcommand{\struck}[1]{\sout{#1}}

% Note du transcripteur, sur l'état matériel du feuillet.
\newcommand{\note}[1]{\par\smallskip{\small\color{gray!60!black}\textit{[#1]}}\par\smallskip}

% Note marginale de Grothendieck, distincte du corps du texte.
\newcommand{\marginal}[1]{\par\smallskip\noindent\hspace{1em}%
  {\small\color{gray!60!black}#1}\par\smallskip}

% --- Titre ------------------------------------------------------------------

\AtBeginDocument{%
  \begin{center}
    {\large\bfseries Fonds Alexandre Grothendieck --- cote n\textsuperscript{o}~\grothFolder}\\[2pt]
    {\itshape\grothFoldertitle}\\[4pt]
    {\small feuillets \grothFirst--\grothLast\ (lot \grothBatch)}\\[2pt]
    {\footnotesize\color{gray!60!black}Transcription établie d'après le fac-similé numérisé
     par l'Université de Montpellier.\\
     Les lectures incertaines sont soulignées, les illisibles notées [\,\ldots\,],
     les ajouts entre crochets.}
  \end{center}
  \vspace{1em}}

\endinput


\folder{11}
\batch{3}
\pages{41}{57}
\dating{[à partir de 1961]}
\watermark{Édition de démonstration}
\foldertitle{Formalisme des correspondances : notes manuscrites (s.d.)}

\begin{document}

\section*{Catégorie des $k$-algèbres graduées bivariantes}

\note{Les titres de section reprennent ses propres mots en tête de chaque
suite ; ceux des pages 43 et 50 sont soulignés ou isolés sur la page. Les mots
qu'il ajoute entre les lignes sont donnés à l'endroit où son renvoi les place.}

\page{41}

Catégorie des $k$-algèbres [graduées] \ill{} bivariantes (anticommutation,
munies d'un degré $d(A)$, nulles en degré $\notin [\uncertain{1}, d(A)]$ ;
\uncertain{on se donne} \ill{} $\kappa_{A} : A \to k$).
\note{Au-dessus de « [graduées] », une insertion de deux mots, reliée par un
trait, ne se lit pas. « anticommutation, » est écrit au-dessus de « munies ».
Le chiffre de l'intervalle se lit $1$ ; l'objet $k$ en degré $0$ de la ligne
suivante demanderait $0$.}

Objet particulier de cette catégorie : $k$ en degré $0$.

Algèbre graduée bivariante jaugée : munie d'un \uncertain{hom.} dans $k$.
Dans certains cas d'hom.\ de telles, $f_{*}$ est déterminé par $f^{*}$.~(*)

Produit tensoriel d'algèbres bivariantes (resp.\ jaugées). Dans certains cas,
c'est un produit dans la catégorie des algèbres bivariantes jaugées.

Algèbre graduée bivariante simpliciale (resp.\ gr.\ biv.\ jaugée)
\struck{\ill{}}. Cas spécial : on a
\[
  A(I' \sqcup I'') \xleftarrow[\ \sim\ ]{(\alpha^{*},\,\beta^{*})} A(I') \otimes A(I'') ,
\]
d'autre part si $\alpha : I \to J$, $\beta : I' \to J'$, donc
$\alpha \sqcup \beta = \gamma : I \sqcup I' \to J \sqcup J'$, on veut que
$(\alpha \sqcup \beta)_{*} : A(I \sqcup I') \to A(J \sqcup J')$ s'identifie à
$\alpha_{*} \otimes \beta_{*}$.
\note{La source et le but de $(\alpha \sqcup \beta)_{*}$ sont récrits l'un sur
l'autre ; on donne la lecture que demandent $\alpha$ et $\beta$.}
On voit que si $A(\Delta^{0}) = A_{0}$ est « de Poincaré » (cas jaugé)
$A(\struck{\ill{}}) \to A(0)$ est une équiv.\ des $A(\cdot)$ avec $A(0)$ de
Poincaré, et des alg.\ de Poincaré. Et $A(\cdot) \mapsto (A(0), \delta)$ est
un foncteur pl.\ fid.\ de la catégorie des alg.\ gr.\ bivariantes jaugées
simpliciales, dans la catégorie des algèbres gr.\ bivariantes jaugées $A$,
munies d'un $\Delta \in [A \otimes A]^{d(A)}$ [il convient de préciser les
axiomes à mettre sur ce $\Delta$].
\note{« gr. » est ajouté au-dessus de « algèbres ». Le $\Delta$ est repassé
deux fois à l'encre forte.}

\marginal{à vérifier ; il faut d'abord \ill{} connaître
$\delta_{*} : A \to A \otimes A$ ; en posant $\Delta = \delta_{*}(1)$}

\textbf{Ex} 1) Si on a une théorie coh.\ \struck{sur la variété}
[satisfaisant Künneth et Poincaré] sur la catégorie $\mathcal{V}$ des variétés
complètes \ill{}, i.e.\ une \uncertain{trivialisation} en dim $d$
\struck{\ill{}}, d'où un foncteur \uncertain{cov.}\ de $\mathcal{V}$ dans la
catégorie $\mathcal{C}_{\mathrm{jaug}}$ des alg.\ gr.\ bivar.\ jaugées sur
$k$, et en transformant par ce foncteur l'objet simplicial canonique
$I \mapsto X^{I}$ défini par $X \in \mathcal{V}$, on trouve une algèbre
bivariante simpliciale bivariante jaugée canonique, d'ailleurs définie aussi
canoniquement, comme dit plus haut, par le composant $H^{*}(X)$ ...
\marginal{\ill{}}
\note{Au-dessus de la fin de « dim $d$ », deux mots remplacent un passage
biffé ; on n'en lit que le premier, \uncertain{resp.}, et un
\uncertain{couple} douteux. Le mot souligné qu'on donne pour
« trivialisation » est lu d'après l'argument : la jauge d'une variété de
dimension $d$. En marge gauche, en face de cette ligne, un mot souligné.}

(*) Attention, les isom.\ $(f^{*}, f_{*})$ ne sont pas néc.\ tels que
$f_{*} = f^{*-1}$ i.e.\ $f_{*}(1) = 1$ ($\Rightarrow$ \struck{\ill{}} que
$\kappa_{A} f^{*} = \kappa_{B}$), ex : $f^{*} = \mathrm{id}$,
$f_{*}(y) = y\sigma$ ($\sigma \in A^{*}$),
$\kappa_{B}(y) = \kappa_{A}(y\sigma^{-1})$ ...

\page{43}

Variantes évidentes avec des variétés topologiques compactes. On peut prendre
pour $k$ un anneau (\uncertain{par exemple}), lorsqu'on a une théorie de
cohomologie à coefficients dans $k$ \struck{qui \ill{}} telle que
\struck{\ill{}} $H^{*}(X)$ satisfait à Poincaré \struck{\ill{}}
(\uncertain{pas seulement} mod torsion ...).
\note{« (par exemple) » est écrit au-dessus de « anneau », « telle que » au-dessus
d'un mot biffé.}

2) Considérons une autre foncteur de $\mathcal{V}$ dans les anneaux bivariants
jaugés, tels les foncteurs $A^{*}(X)$ (anneau de Chow, \struck{gradué} jaugé
par le degré des \uncertain{0-}cycles de degré \uncertain{max.}), ou un
foncteur quotient (tels les groupes de cycles mod équival.\ numérique,
homologique, $\tau$-algébrique, Picard, Albanese ...). On trouve un anneau
bivariant jaugé simplicial, qui n'a plus aucune tendance à être spécial, et
dont les composantes n'ont pas trop envie d'être de Poincaré ...
\note{Le $\tau$ est ajouté entre les lignes, avec un renvoi devant
« algébriques ».}

On peut tensoriser par $\mathbf{Q}$ ....

3) On peut prendre $K(X)$ [jaugé \uncertain{par} $\chi$], considéré comme
gradué de degré $0$ ....

\section*{Formalisme des correspondances}

Formalisme des correspondances. Soit $C \in A(I \sqcup I')$. On en déduit une
application
\[
  x \longmapsto C(x) = C.x \quad : \quad A(I) \to A(I')
\]
comme le composé (où $\alpha : I \to I \sqcup I'$, $\alpha' : I' \to I \sqcup I'$)
\[
  A(I) \xrightarrow{\ \alpha^{*}\ } A(I \sqcup I') \xrightarrow{\ y \mapsto Cy\ }
  A(I \sqcup I') \xrightarrow{\ \alpha'_{*}\ } A(I') .
\]
\marginal{il y a un choix : prendre $Cy$ ou $yC$ !}
\note{Sous la flèche $\mapsto$ de la première formule, un petit signe ne se
lit pas ; au-dessus de $\alpha^{*}$, une flèche vers la gauche est surchargée.}

Ces opérations sont connues quand on connaît, en plus des lois
d'\emph{anneaux} des $A(\cdot)$, les $\alpha^{*}$ et $\alpha_{*}$ pour
$\alpha : I \hookrightarrow J$ une injection [il suffit même
$\operatorname{card}(J) = \operatorname{card}(I) + 1$, i.e.\ $\alpha$ une
« opération face »].

\page{45}

\uncertain{Posons} $A(I', I)$ au lieu de $A(I \sqcup I')$.
\struck{Corresp.}\ ($A(I \sqcup I')$ \struck{\ill{}} corresp.\ de $I$ à $I'$)
[On va en définir une opération naturelle bilinéaire
\[
  A(I'', I') \times A(I', I) \longrightarrow A(I'', I)
\]
notée $(v, u) \longmapsto v \circ u$, et définie par
\[
  v \circ u = p_{31*}\bigl(p_{21}^{*}(u)\, p_{32}^{*}(v)\bigr) .
\]
\note{Dans $p_{32}^{*}(\cdot)$, ici et deux fois au bas de la page, la lettre
est tracée comme un $w$ ; la formule demande $v$, et on lit $v$.}

\begin{tikzcd}
  & X^{I \sqcup I'} \arrow[dl, "\alpha"'] \arrow[dr, "\alpha'"] & \\
  X^{I} & & X^{I'}
\end{tikzcd}

Ceci posé, si $\widetilde{u} : A(I) \to A(I')$, \struck{\ill{}} ..., on
vérifie la formule
\[
  \widetilde{v \circ u} = \widetilde{v} \circ \widetilde{u} .
\]

\begin{tikzcd}[column sep=small, row sep=small, nodes={font=\scriptsize}]
  & & & X^{I \sqcup I' \sqcup I''} \arrow[dll, "p_{21}"'] \arrow[d, "p_{31}"] \arrow[drr, "p_{32}"] & & & \\
  & X^{I \sqcup I'} \arrow[dl, no head, "\gamma"'] \arrow[dr, "\gamma'"] & & X^{I \sqcup I''} \arrow[ddl, no head, "\beta"'] \arrow[ddr, no head, "\beta''"] & & X^{I' \sqcup I''} \arrow[dl, no head, "\alpha'"'] \arrow[dr, "\alpha''"] & \\
  X^{I} & & X^{I'} & & X^{I'} & & X^{I''} \\
  & & X^{I} & & X^{I''} & &
\end{tikzcd}

\note{Les traits $\gamma$, $\beta$, $\beta''$ et $\alpha'$ sont tracés sans
pointe sur la page, les quatre autres avec ; on les donne tels.}

En effet, on a
\begin{align*}
  (\widetilde{v \circ u})(x)
  &\overset{\text{déf.\ de }\sim}{=} \struck{\ill{}}\; \beta''_{*}\bigl[(v \circ u)\, \beta^{*}(x)\bigr] \\
  &\overset{\text{déf.\ de }v \circ u}{=} \beta''_{*}\bigl[\underbrace{p_{31*}\bigl(p_{21}^{*}(u)\, p_{32}^{*}(v)\bigr) \beta^{*}(x)}_{=\ \text{formule de proj.}}\bigr] \\
  &= \beta''_{*}\, p_{31*}\bigl(p_{21}^{*}(u)\, p_{32}^{*}(v)\, \underbrace{p_{31}^{*} \beta^{*}(x)}_{p_{1}^{*}(x)}\bigr) \\
  &\overset{\text{transitivité}}{=} p_{3*}\bigl(p_{21}^{*}(u)\, p_{32}^{*}(v)\, p_{1}^{*}(x)\bigr) \\
  &\struck{\overset{\text{projection}}{=} p_{3*}\bigl(p_{21}^{*}(u)\, p_{32}^{*}(v)\bigr) x}
\end{align*}
\note{Au troisième membre, le $\beta''_{*}$ de tête n'est pas récrit sur la
page ; on le rétablit pour que la chaîne se lise. La dernière ligne est
biffée.}

\begin{tikzcd}
  & X^{I \sqcup I' \sqcup I''} \arrow[dl, "p_{21}"'] \arrow[dr, "p_{32}"] & & \\
  X^{I \sqcup I'} \arrow[dr, "\gamma'"'] & & X^{I' \sqcup I''} \arrow[dl, "\alpha'"] \arrow[dr, no head, "\alpha''"] & \\
  & X^{I'} & & X^{I''}
\end{tikzcd}

\note{L'exposant du sommet est surchargé ; un trait fin sépare ce carré du
calcul qui précède.}

\begin{align*}
  \widetilde{v}\bigl(\widetilde{u}(x)\bigr)
  &\overset{\text{déf.\ }\widetilde{v}}{=} \struck{\ill{}}\; \alpha''_{*}\bigl(v\, \alpha'^{*}(\widetilde{u}(x))\bigr) \\
  &\overset{\text{déf.\ }\widetilde{u}}{=} \alpha''_{*}\bigl(v\, \underbrace{\alpha'^{*}\bigl(\gamma'_{*}(u\, \gamma^{*}(x))\bigr)}_{\text{échange}}\bigr)
\end{align*}
\marginal{« axiome cartésien » oublié !}
\[
  \alpha'^{*}\gamma'_{*}\bigl(u\,\gamma^{*}(x)\bigr) = \struck{\ill{}}\;
  p_{32*}\bigl(p_{21}^{*}(u\, \gamma^{*}(x))\bigr)
  = p_{32*}\bigl(p_{21}^{*}(u)\, \underbrace{p_{21}^{*} \gamma^{*}(x)}_{p_{1}^{*}(x)}\bigr)
\]
\note{La bulle « axiome cartésien oublié ! » est tirée vers le signe $=$ de
l'échange.}
\begin{align*}
  &= \alpha''_{*}\bigl(v\, \underbrace{p_{32*}\bigl(p_{21}^{*}(u)\, p_{1}^{*}(x)\bigr)}_{\text{formule de projection, mod signes}}\bigr) \\
  &= \underbrace{\alpha''_{*}\, p_{32*}}\bigl(p_{32}^{*}(v)\, p_{21}^{*}(u)\, p_{1}^{*}(x)\bigr) \\
  &= p_{3*}\bigl(p_{32}^{*}(v)\, p_{21}^{*}(u)\, p_{1}^{*}(x)\bigr)
\end{align*}
\marginal{OK mod signes ; faire attention}
\note{Un $= \alpha$ biffé précède l'avant-dernière ligne.}

\page{47}

On va définir aussi
\[
  \Delta_{I} \in A(I \sqcup I)
\]
qui \uncertain{sera} tel que
\[
  \widetilde{\delta}_{I}(x) = x \quad (x \in A(I)) \qquad \text{i.e.}\quad
  \widetilde{\delta}_{I} = \mathrm{id}_{A(I)} .
\]
\note{Sous le tilde, la lettre est un $\delta$ minuscule, là où la ligne
précédente pose $\Delta_{I}$.}

Pour \uncertain{le voir}, on \struck{définit} considère
$\delta : I \sqcup I \to I$ d'où $\delta_{*} : A(I) \to A(I \sqcup I)$, on pose
\[
  \Delta_{I} = \delta_{*}(1) .
\]

On a (écrivant $\Delta$ au lieu de $\Delta_{I}$)

\begin{tikzcd}
  & X^{I} \arrow[d, "\delta"] & \\
  & X^{I \sqcup I} \arrow[dl, "p_{1}"'] \arrow[dr, "p_{2}"] & \\
  X^{I} & & X^{I}
\end{tikzcd}

\begin{align*}
  \widetilde{\Delta}(x)
  &\overset{\text{déf.\ de }\sim}{=} p_{2*}\bigl(\Delta\, p_{1}^{*}(x)\bigr) \\
  &= p_{2*}\bigl(\underbrace{\delta_{*}(1)\, p_{1}^{*}(x)}_{\delta_{*}(\delta^{*} p_{1}^{*}(x)) \,=\, \delta_{*}(x)}\bigr) \\
  &= p_{2*}\, \delta_{*}(x) = x
\end{align*}
\note{Devant le $\delta_{*}$ de l'accolade, un $p$ est surchargé.}

Il faut encore \uncertain{définir} l'associativité de l'opération $\circ$, et
le fait que les $\delta_{I}$ jouent le rôle d'unités. On voit que sous ces
conditions,
\note{La page s'arrête sur cette virgule.}

\section*{Sous-algèbre tautologique de $H^{*}(X^{n})$}

\page{50}

Sous-algèbre tautologique de $H^{*}(X^{n})$.

(1) $H^{*}(X)$ \struck{degrés}

\begin{itemize}
  \item \emph{générateurs} : $y \in H^{n}(X)$
  \item \emph{relations} : $y^{2} = 0$
  \item \emph{base} : deg.\ $0$ : $1$ ; deg.\ $n$ : $y$
\end{itemize}

(2) $H^{*}(X \times X)$

\begin{itemize}
  \item \emph{générateurs} : $y \otimes 1 = y_{1}$, $1 \otimes y = y_{2}$,
  $\delta \in H^{n}(X)$
  \item \emph{relations} : $y_{1}^{2} = y_{2}^{2} = 0$,
  $y_{1} y_{2} = y_{1} \delta = y_{2} \delta$ ; $\delta^{2} = \chi\, y_{1} y_{2}$
  ($\chi \in \mathbf{Z}$ la car.\ d'EP)
  \item \emph{base} : deg $0$ : $1$ ; deg $n$ : $y_{1}, y_{2}, \delta$ ;
  deg $2n$ : $y_{1} y_{2} y_{3}$
\end{itemize}
\note{« $y_{1} y_{2} y_{3}$ » est écrit tel, sic : sur $X \times X$ il n'y a
que $y_{1}$, $y_{2}$, et le degré $2n$ demande $y_{1} y_{2}$.
« $\delta \in H^{n}(X)$ » est écrit tel ; le $\delta$ vit dans
$H^{n}(X \times X)$. Les numéros (1), (2), (3) et ($k$) sont cerclés sur la page. Après $y_{1}$, un signe est biffé.}

(3) $H^{*}(X \times X \times X)$

\begin{itemize}
  \item \emph{générateurs} : deg.\ $n$ : $y \otimes 1 \otimes 1 = y_{1}$,
  $1 \otimes y \otimes 1 = y_{2}$, $1 \otimes 1 \otimes y = y_{3}$,
  $\delta_{12} = \delta \otimes 1$, $\delta_{23}$, $\delta_{31}$ ;
  deg $2n$ : $\delta$
  \item \emph{relations} :
  \[
  \left\{
  \begin{array}{l}
    y_{1}^{2} = y_{2}^{2} = y_{3}^{2} = 0 ,\\
    \delta_{12}^{2} = \chi\, y_{1} y_{2} ,\ \delta_{23}^{2} = \chi\, y_{2} y_{3} ,\quad
    \delta_{31}^{2} = \chi\, y_{3} y_{1} ,\quad
    \delta_{12} \delta_{23} = \delta_{12} \delta_{31} = \delta_{23} \delta_{31} = \delta \\
    \delta_{12} y_{1} = \delta_{12} y_{2} = y_{1} y_{2} ,\quad
    \delta_{23} y_{2} = \delta_{23} y_{3} = y_{2} y_{3} ,\quad
    \delta_{31} y_{3} = \delta_{31} y_{1} = y_{3} y_{1} ,\\
    \delta y_{1} = \delta y_{2} = \delta y_{3} = y_{1} y_{2} y_{3} \\
    \delta^{2} = 0 \\
    \delta_{12} \delta = \delta_{23} \delta = \delta_{31} \delta = \chi\, y_{1} y_{2} y_{3}
  \end{array}
  \right.
  \]
  \item \emph{base} : deg $0$ : $1$ ;
  deg $n$ : $y_{1}, y_{2}, y_{3}, \delta_{12}, \delta_{23}, \delta_{31}$ ;
  deg $2n$ : $y_{1} y_{2}, y_{2} y_{3}, y_{3} y_{1}, \delta_{12} y_{3},
  \delta_{23} y_{1}, \delta_{31} y_{2}, \delta$ ;
  deg $3n$ : $y_{1} y_{2} y_{3}$
\end{itemize}

\page{52}

($k$) $H^{*}(X^{k})$ \struck{générateurs}

\emph{générateurs}
\begin{itemize}
  \item \emph{degré $n$} : les $k$ transformés de
  $y \otimes \underbrace{1 \otimes \cdots \otimes 1}_{k-1} = \overline{y}$ ;
  les $\binom{k}{2}$ transformés de
  $\delta_{1} \otimes \underbrace{1 \otimes \cdots \otimes 1}_{k-2} = \overline{\delta}_{1}$
  \item \emph{degré $2n$} : les $\binom{k}{3}$ transformés de
  $\delta_{2} \otimes \underbrace{1 \otimes \cdots \otimes 1}_{k-3} = \overline{\delta}_{2}$
  \item ---
  \item \emph{degré $rn$} : les $\binom{k}{r+1}$ transformés de
  $\delta_{r} \otimes \underbrace{1 \otimes \cdots \otimes 1}_{k-1-r} = \overline{\delta}_{r}$
  \item ---
  \item \emph{degré $(k-1)n$} : $\delta_{k-1} = \overline{\delta}_{k-1}$
\end{itemize}
\note{Le $r$ de « degré $rn$ » est écrit à l'encre forte sur un $k$. Sur
$\delta_{1}$ un petit signe en exposant ne se lit pas.}

\struck{relations}

\emph{relations}
\note{La page s'arrête sur ce mot.}

\page{54}

\[
  X^{k-i} \longrightarrow X^{k}
  \qquad
  \xleftarrow{\ \ill{}\ } [1, k'] \subset [1, k]
\]

\[
  X^{I}
\]

$I = I' \sqcup I''$ \qquad $R$ relation d'équivalence sur $I'$

$I$ avec une division en \uncertain{un ens.}\ de parties \struck{non vides},
comprenant la partie vide, et le choix d'une de ses parties

\[
  (R, J) \quad
  \left\{
  \begin{array}{l}
    R \subset \mathfrak{P}(I) \text{ tel que } \struck{\ill{}}
    \left\{
    \begin{array}{l}
      \text{a)}\ \bigcup_{A \in R} A = I \\
      \text{b)}\ A, A' \in R,\ A \neq A' \Rightarrow A \cap A' = \emptyset \\
      \text{c)}\ \emptyset \in R
    \end{array}
    \right. \\[1ex]
    J \in R
  \end{array}
  \right.
\]
\note{Dans a) l'indice $A \in R$ est écrit sur $J \in R$ ; dans b) les $A$
sont repassés sur d'autres lettres.}

\[
  \varepsilon_{R,J}\, \varepsilon_{R',J'} =
  \left\{
  \begin{array}{l}
    0 \quad \text{si } J \cap J' \neq \emptyset \\
    \phantom{0}
  \end{array}
  \right.
\]

Supposons $J \cap J' = \emptyset$
\[
  Z_{R,J} = \left\{ (x_{i})_{i \in I} \text{ tel que }
  \left\{
  \begin{array}{ll}
    x_{i} = a & \text{si } i \in J \\
    x_{i} = x_{j} & \text{si } \exists A \in R, \text{ t.q. } i, j \in A
  \end{array}
  \right.
  \right\}
\]
\[
  Z^{a}_{R,J} \cap Z^{a'}_{R',J'} \underset{\text{du pt de vue ens.}}{=} Z_{R'',J''}
\]
$R''$ est la partition, engendrée par $R, R'$

$\widetilde{J}$ la \uncertain{classe} \struck{\ill{}} engendrée par $J$

$\widetilde{J}'$ --- engendrée par $J'$
\note{Le deuxième « engendrée » est souligné sous la ligne, le premier biffé
dans la ligne de $\widetilde{J}$.}

\[
  \struck{\widetilde{J} \neq \widetilde{J}'}
\]
\[
  \left|
  \begin{array}{l}
    \widetilde{J} = \widetilde{J}' \neq \emptyset \Rightarrow
    \varepsilon_{R,J}\, \varepsilon_{R',J'} = 0 \\[1ex]
    \widetilde{J} \neq \widetilde{J}' \text{ ou } \widetilde{J} = \widetilde{J}' = \emptyset
    \quad \text{alors} \quad
    \varepsilon_{R,J}\, \varepsilon_{R',J'} =
    \Bigl(\prod_{i \in \widetilde{J} \cup \widetilde{J}'} y_{i}\Bigr)
    \prod_{\substack{A'' \in R'' \\ A'' \neq \widetilde{J} \\ A'' \neq \widetilde{J}'}}
    \varepsilon_{R_{A''}}\, \varepsilon_{R'_{A''}}
  \end{array}
  \right.
\]
\note{Les exposants et les indices des deux derniers $\varepsilon$ sont
surchargés au point de ne plus se lire ; on ne donne que ce qui reste
visible, $R$, $R'$ et $A''$.}

On est ramené au cas où $J = \emptyset$, $J' = \emptyset$, et la relation
d'équivalence engendrée par $R, R'$ est la relation grossière.

Alors, \uncertain{ensemblistement}, on obtient la diagonale de $X^{I}$

\page{56}

\[
  1 \leq \operatorname{card} R = s \leq k \quad \text{donc} \quad
  \dim Z_{R} = sn \qquad \operatorname{codim} Z_{R} = kn - sn = (k-s)n
\]
\[
  \operatorname{card} R' = s' \qquad \operatorname{codim} Z_{R'} = (k-s')n
\]
\[
  \operatorname{card} R'' = 1 \qquad \operatorname{codim} Z_{R''} = (k-1)n
\]
\[
  (k-1)n \leq (k-s)n + (k-s')n \quad \text{i.e.} \quad \struck{\ill{}}
\]
\[
  \boxed{s + s' \leq k+1}
\]

\marginal{à vérifier\quad OK}
Cas d'égalité \quad $s + s' = k+1$ : on trouve la diagonale $\delta_{k}$.

\marginal{à vérifier}
Cas d'inégalité \quad $k+1 > s + s'$ i.e.\ $k \geq s + s'$
\begin{quote}
  \struck{donc il faut avoir} \struck{$k+1 > s+s'$}
  1) $k > s + s'$ alors $\varepsilon_{R}\, \varepsilon_{R'} = 0$ pour raison
  de degrés

  2) $k = s + s'$, alors $\varepsilon_{R}\, \varepsilon_{R'} =
  \struck{\chi}\; \uncertain{y}_{X^{k}}$ \quad ?
\end{quote}
\note{« $s+s' = k+1$ » est souligné d'un trait ondulé. Dans « $k \geq s+s'$ »
le signe est écrit sur un autre. Au cas 2), le point d'interrogation est
tiré sous la classe par un trait.}

\begin{tikzcd}
  X^{s} \arrow[r, hook] & X^{k} \\
  & X^{s'} \arrow[u, hook]
\end{tikzcd}

$s + s' = k$ ou $k+1$

\struck{$k \leq k - s + k - s'$}

\begin{tikzcd}
  V^{s} \arrow[r] & V^{k} \\
  V \arrow[u] \arrow[r] & V^{s'} \arrow[u]
\end{tikzcd}

\begin{tikzcd}
  & U \arrow[dr, no head, "p"] & \\
  U \cap V \arrow[ur, no head] \arrow[dr, no head] & & W \\
  & V \arrow[ur, no head, "q"'] &
\end{tikzcd}

\begin{tikzcd}
  & U' \arrow[dr, no head, "p-r"] & & \\
  U' \cap V' \arrow[ur, no head, "q-r"] \arrow[dr, no head, "p-r"'] & & U' + V' \arrow[r, no head, "r"] & W' \\
  & V' \arrow[ur, no head, "q-r"'] & &
\end{tikzcd}

\note{Dans le premier losange, les traits portent à un bout la boucle d'une
inclusion, sans pointe ; à gauche sont écrits $p+q-r$ et, plus bas, $r$. Dans
le second, le $q-r$ de $U' \cap V'$ à $U'$ est écrit sur un $p$.}

\begin{itemize}
  \item a) Stable par $\boxtimes$ : clair
  \item b) Stable par images inverses par projections simpliciales : clair
  \item c) Stable par images \emph{directes} par projections :
  $i \in I$ \uncertain{définit} \ill{} projections $X^{I} \to X^{I - \{i\}}$
  \item \quad a) $i \in J$ : clair
  \item \quad b) $i \notin J$ : on est ramené au cas de
  $\delta_{I} \in H^{*}(X^{I}) \to H^{*}(X^{I - \{i\}})$
  \struck{\ill{}} on trouve $\delta_{I}$
  \item d) \struck{\ill{}}
\end{itemize}
\note{Le d) et les trois lignes qui le suivent, dont une encadrée, sont biffés
et surchargés ; on y devine « Stable par », « les classes », sans plus. Dans
« on trouve $\delta_{I}$ », l'indice porte un trait.}

\[
  X^{p} \to X^{q} \qquad X^{p} \times X \to X^{q} \times X \qquad
  \Delta_{p} \leftarrow \Delta_{q} \qquad
  \Delta_{p} \sqcup 1 \leftarrow \Delta_{q} \sqcup 1
\]

Stable par comp.

\[
  \boxed{\text{engendrée, via } \boxtimes, \text{ par } \struck{\ill{}}\ y, \text{ et les } \delta_{I}}
\]

\[
  (\delta_{2} \boxtimes \delta_{2})
\]

\end{document}
